\section{Conclusion}
\label{sec:conclusion}

We began by observing that every major LLM pipeline uses data curation, yet optimal filtering parameters remain undiscovered. Our work demonstrates that this gap can be addressed through systematic dose-response analysis, treating curation parameters as continuous variables rather than discrete configuration choices.

\subsection{Summary}

In this work, we addressed the lack of controlled ablation studies for LLM data curation by applying fixed-token experimental design to map parameter effect surfaces. Our key insight---that perplexity filtering exhibits a concave dose-response relationship due to the quality-diversity tradeoff---enables principled threshold selection.

Our main contributions are:

\begin{enumerate}
    \item \textbf{Existence of non-monotonic dose-response.} We demonstrated that the relationship between perplexity thresholds and benchmark performance is quadratic ($R^2=0.985$), not linear. This validates the quality-diversity tradeoff as an empirical phenomenon and provides the first controlled evidence that intermediate thresholds outperform both extremes.

    \item \textbf{Mechanistic validation through convergence analysis.} We showed that noise dilution accounts for the left side of the curve (unfiltered training shows 40\% higher convergence AUC) while diversity loss explains the right side (p90 filtering degrades final performance despite reasonable convergence). This mechanistic understanding supports principled rather than heuristic parameter selection.

    \item \textbf{Practical guidance with scale transfer.} We demonstrated 1.32\% improvement over RedPajama defaults and showed that optimal thresholds transfer across model scales with ratio 0.85, enabling efficient optimization at smaller scales with predictable extrapolation.
\end{enumerate}

\subsection{Future Directions}

This work opens several promising directions grounded in our experimental findings.

\textbf{From untested alternative explanations:} Our experiments used only GPT-2 architecture. The non-monotonic pattern may be architecture-specific rather than universal. Future work should replicate the dose-response sweep with Llama-style architectures to test generalization.

\textbf{From unverified assumptions:} We assumed that KenLM 5-gram perplexity is an appropriate quality signal. Alternative signals---classifier-based quality scores, BERT perplexity, or semantic filtering---may yield different optimal thresholds or sharper dose-response curves.

\textbf{From scope extensions:} Our experiments studied perplexity and deduplication independently. Joint optimization across both dimensions may reveal interaction effects and a more efficient Pareto frontier.

\subsection{Closing}

Data curation has long been treated as a pipeline preprocessing step---important but unoptimized. Our findings suggest it deserves the same systematic attention as architectural choices or training hyperparameters. The optimal perplexity threshold of approximately p50 provides a validated starting point, but more importantly, the dose-response methodology offers a framework for continued optimization as corpora and architectures evolve.

We hope this work encourages the community to treat curation parameters as first-class optimization targets, transforming data preparation from craft to science.
